\documentclass[11pt,a4paper]{article}

% --- Packages ---
\usepackage[margin=.75in]{geometry}
\usepackage[utf8]{inputenc}
\usepackage[T1]{fontenc}
\usepackage{lmodern} % Fix for font not found error
\usepackage{titlesec}
\usepackage{enumitem}
\usepackage{hyperref}
\usepackage{xcolor}
\usepackage{fancyhdr}
\usepackage{comment}
\usepackage{cleveref}
\usepackage[bottom]{footmisc}
\usepackage{graphicx}
\usepackage{longtable}
\usepackage{booktabs}
\usepackage{array}
\usepackage{calc}
\usepackage{etoolbox}

% Standard Pandoc setup
\providecommand{\tightlist}{%
  \setlength{\itemsep}{0pt}\setlength{\parskip}{0pt}}

\crefformat{section}{\S#2#1#3}
\crefformat{subsection}{\S#2#1#3}
\crefformat{subsubsection}{\S#2#1#3}
\crefformat{figure}{#2Figure~#1#3}
\crefformat{footnote}{#2\footnotemark[#1]#3}

\linespread{0.9} % Riduce lo spazio tra le linee

% --- Colors ---
\definecolor{primary}{RGB}{0, 0, 0}
\definecolor{links}{HTML}{0000EE}
\definecolor{blue}{HTML}{26428B}

\hypersetup{
    colorlinks=true,
    linkcolor=links,
    urlcolor=links,
    citecolor=links,
}

% --- Section Formatting ---
\titleformat{\section}
  {\normalfont\Large\bfseries\color{blue}}{\thesection}{1em}{}
\titleformat{\subsection}
{\normalfont\large\bfseries\color{blue}}{\thesubsection}{1em}{}
\titleformat{\subsubsection}
{\normalfont\bfseries\color{blue}}{\thesubsubsection}{1em}{}

% --- Header and Footer ---
\pagestyle{fancy}
\fancyhead[L]{\small Giuseppe Capasso}
\fancyhead[R]{\small devnet-workshop}
\fancyfoot[C]{\thepage}
\renewcommand{\headrulewidth}{0.4pt}

% --- Custom Title Command ---
\newcommand{\proposaltitle}[1]{
    \begin{center}
        \vspace*{0.1cm}
        {\LARGE \bfseries \color{blue} #1} \\[1.5ex]
        {\large \textbf{Giuseppe Capasso}} \\[1ex]
                \small{
            \vspace{0.1cm}
            \textbf{Materia:} devnet-workshop\\
        }
            \end{center}
    \vspace{0.3cm}
    \hrule height 0.5pt
    \vspace{0.5cm}
}

\begin{document}

\proposaltitle{Docker Hello world}

\section{Docker Hello world}\label{docker-hello-world}

\subsection{Installing Docker}\label{installing-docker}

Visit the web \href{https://docs.docker.com/desktop/}{page} and click on
the correct link based on your OS. Next, to the section matching your
environment.

Docker desktop homepage

Docker desktop install

\subsubsection{POV: you are on Windows}\label{pov-you-are-on-windows}

To run Docker on Windows, you need some extra step.

You need to install WSL2 (Windows subsystem for linux). Other
requirements listed in the page should be automatically satisfied if you
have successfully installed a Linux virtual machine.

Windows docker desktop

\paragraph{Installing WSL}\label{installing-wsl}

To install WSL, visit the
\href{https://learn.microsoft.com/en-us/windows/wsl/install?source=docs}{docs}.

Open a \texttt{Powershell\ (Administrator\ prompt)} and run the command
\texttt{wsl\ -\/-install}.
\includegraphics[width=5.83333in,height=2.89821in]{media/rId31.png}

Now you should have an Ubuntu CLI integrated in Windows!

\paragraph{Installing Windows
Terminal}\label{installing-windows-terminal}

Windows terminal is an utility used to create and manage terminal like
in Linux. You can install it from the
\href{https://learn.microsoft.com/en-us/windows/terminal/install}{official
documentation}.

Windows docker desktop

\subsubsection{POV: you are on MAC (intel/arm
chip)}\label{pov-you-are-on-mac-intelarm-chip}

You should be fine: Docker is supported on recent version of macOS.

Mac docker desktop

Just follow the instruction and you should have a working Docker
installation.

Mac docker desktop installation

\subsubsection{POV: you are on Linux}\label{pov-you-are-on-linux}

Docker is natively supported by Linux. Before installing Docker desktop
we need to install the Docker engine. For example, if you are on a
Debian based distro, you can find the
\href{https://docs.docker.com/engine/install/ubuntu/\#install-using-the-repository}{docs}

Linux docker engine install

This script sets up your repository, now you can install by running the
command:

\begin{verbatim}
sudo apt-get install docker-ce docker-ce-cli containerd.io docker-buildx-plugin docker-compose-plugin
\end{verbatim}

Now, if you want you can install Docker Desktop but it is not necessary
for this lab.

\paragraph{Permission issues}\label{permission-issues}

If you are on Linux, you can run containers with \texttt{sudo}. If you
want to execute \texttt{docker} as a normal user you can put yourself in
the \texttt{docker\ group} running the command:

\begin{verbatim}
sudo usermod -aG docker <your username>
\end{verbatim}

Restart your machine and now you are ready to test the installation.

\subsection{Testing your docker
installation}\label{testing-your-docker-installation}

Every one should have Docker installation. To check if everything worked
open a terminal (Windows terminal for Windows users) and run the
command:

\begin{verbatim}
docker run hello-world
\end{verbatim}

You should get an output like this:

\begin{verbatim}
Unable to find image 'hello-world:latest' locally
latest: Pulling from library/hello-world
c1ec31eb5944: Pull complete
Digest: sha256:91bc16c380fe750bcab6a4fd29c55940a7967379663693ec9f4749d3878cd939
Status: Downloaded newer image for hello-world:latest

Hello from Docker!
This message shows that your installation appears to be working correctly.

To generate this message, Docker took the following steps:
 1. The Docker client contacted the Docker daemon.
 2. The Docker daemon pulled the "hello-world" image from the Docker Hub.
    (amd64)
 3. The Docker daemon created a new container from that image which runs the
    executable that produces the output you are currently reading.
 4. The Docker daemon streamed that output to the Docker client, which sent it
    to your terminal.

To try something more ambitious, you can run an Ubuntu container with:
 $ docker run -it ubuntu bash

Share images, automate workflows, and more with a free Docker ID:
 https://hub.docker.com/

For more examples and ideas, visit:
 https://docs.docker.com/get-started/
\end{verbatim}

\subsection{Your first Docker
application}\label{your-first-docker-application}

To test things out, we can deploy a simple python application.

The application will print out information about the operating system
and will ask for a username. The code should look like this.

\begin{verbatim}
import platform
import sys

print("Hello, you are executing me from", platform.machine(), "processor")
print("Platform information:", platform.platform())
print("OS:", platform.system())
print("Python:", sys.version)
name = input("Give me your name: \n")

print("Hello,", name)
\end{verbatim}

Now we need a Dockerfile. For example we decide to use this one:

\begin{verbatim}
FROM python:3.10.14-bookworm

WORKDIR /app
COPY main.py .

CMD [ "python", "/app/main.py" ]
\end{verbatim}

Now we can build the image with the command:

\begin{verbatim}
docker build -t dtlab-docker .
\end{verbatim}

And we execute it

\begin{verbatim}
docker run -i dtlab-docker
\end{verbatim}

\textbf{STOP and THINK}: what does the \texttt{-i} flag do?

\end{document}
